
\documentclass{article}
\usepackage{colortbl}
\usepackage{makecell}
\usepackage{multirow}
\usepackage{supertabular}

\begin{document}

\newcounter{utterance}

\centering \large Interaction Transcript for game `referencegame', experiment `random\_grids', episode 14 with kimi{-}k2.5{-}without{-}reasoning{-}t1.0.
\vspace{24pt}

{ \footnotesize  \setcounter{utterance}{1}
\setlength{\tabcolsep}{0pt}
\begin{supertabular}{c@{$\;$}|p{.15\linewidth}@{}p{.15\linewidth}p{.15\linewidth}p{.15\linewidth}p{.15\linewidth}p{.15\linewidth}}
    \# & $\;$A & \multicolumn{4}{c}{Game Master} & $\;\:$B\\
    \hline

    \theutterance \stepcounter{utterance}  
    & & \multicolumn{4}{p{0.6\linewidth}}{
        \cellcolor[rgb]{0.9,0.9,0.9}{
            \makecell[{{p{\linewidth}}}]{
                \texttt{\tiny{[P1$\langle$GM]}}
                \texttt{You are given three grids, where each of them is 5 by 5 in size.} \\
\texttt{Grids have empty cells marked with "▢" and filled cells marked with "X".} \\
\texttt{Your task is to generate a referring expression that best describes the target grid while distinguishing it from the two other distractor grids.} \\
\texttt{The first grid is the target grid, and the following two grids are the distractors.} \\
\\ 
\texttt{Target grid:} \\
\\ 
\texttt{□ X X □ X} \\
\texttt{X □ X X □} \\
\texttt{□ X □ X X} \\
\texttt{X □ X □ X} \\
\texttt{□ X X □ □} \\
\\ 
\texttt{Distractor grid 1:} \\
\\ 
\texttt{□ X □ X □} \\
\texttt{X □ X □ X} \\
\texttt{□ X □ X □} \\
\texttt{X □ X □ X} \\
\texttt{□ X □ X □} \\
\\ 
\texttt{Distractor grid 2:} \\
\\ 
\texttt{□ □ X X □} \\
\texttt{X □ □ X X} \\
\texttt{□ X X □ X} \\
\texttt{X X □ □ X} \\
\texttt{□ X X □ □} \\
\\ 
\texttt{Instruction: Describe the target grid.} \\
\texttt{Generate the referring expression starting with the tag "Expression: " for the given target grid. Omit any other text.} \\
            }
        }
    }
    & \\ \\

    \theutterance \stepcounter{utterance}  
    & \multicolumn{4}{p{0.6\linewidth}}{
        \cellcolor[rgb]{0.9,0.9,0.9}{
            \makecell[{{p{\linewidth}}}]{
                \texttt{\tiny{[P1$\rangle$GM]}}
                \texttt{Expression: The grid with X's in the top{-}right corner and bottom{-}right area, featuring a diagonal pattern of X's from top{-}left to bottom{-}right that is broken by filled cells in positions (1,2), (2,3), (3,4), and (4,5), with empty cells at the four corners of the bottom{-}right 2x2 subgrid.} \\
            }
        }
    }
    & & \\ \\

    \theutterance \stepcounter{utterance}  
    & & & \multicolumn{2}{p{0.3\linewidth}}{
        \cellcolor[rgb]{0.9,0.9,0.9}{
            \makecell[{{p{\linewidth}}}]{
                \texttt{\tiny{[GM$|$GM]}}
                \texttt{Expression: The grid with X's in the top{-}right corner and bottom{-}right area, featuring a diagonal pattern of X's from top{-}left to bottom{-}right that is broken by filled cells in positions (1,2), (2,3), (3,4), and (4,5), with empty cells at the four corners of the bottom{-}right 2x2 subgrid.} \\
            }
        }
    }
    & & \\ \\

    \theutterance \stepcounter{utterance}  
    & & \multicolumn{4}{p{0.6\linewidth}}{
        \cellcolor[rgb]{0.9,0.9,0.9}{
            \makecell[{{p{\linewidth}}}]{
                \texttt{\tiny{[GM$\rangle$P2]}}
                \texttt{You are given three grids, where each of them is 5 by 5 in size.} \\
\texttt{Grids have empty cells marked with "▢" and filled cells marked with "X".} \\
\texttt{You are also given a referring expression that describes one of the given grids.} \\
\texttt{Your task is to select the grid that matches the given referring expression.} \\
\texttt{Generate only the number (in text) of the grid that matches the given expression by selecting first, second, or third.} \\
\\ 
\texttt{First:} \\
\\ 
\texttt{□ □ X X □} \\
\texttt{X □ □ X X} \\
\texttt{□ X X □ X} \\
\texttt{X X □ □ X} \\
\texttt{□ X X □ □} \\
\\ 
\texttt{Second:} \\
\\ 
\texttt{□ X □ X □} \\
\texttt{X □ X □ X} \\
\texttt{□ X □ X □} \\
\texttt{X □ X □ X} \\
\texttt{□ X □ X □} \\
\\ 
\texttt{Third:} \\
\\ 
\texttt{□ X X □ X} \\
\texttt{X □ X X □} \\
\texttt{□ X □ X X} \\
\texttt{X □ X □ X} \\
\texttt{□ X X □ □} \\
\\ 
\texttt{Expression: The grid with X's in the top{-}right corner and bottom{-}right area, featuring a diagonal pattern of X's from top{-}left to bottom{-}right that is broken by filled cells in positions (1,2), (2,3), (3,4), and (4,5), with empty cells at the four corners of the bottom{-}right 2x2 subgrid.} \\
\texttt{Question: Which grid does the expression refer to?} \\
\texttt{Start with the tag "Answer: ", followed by your selection. Omit any other text.} \\
            }
        }
    }
    & \\ \\

    \theutterance \stepcounter{utterance}  
    & & & \multicolumn{4}{p{0.6\linewidth}}{
        \cellcolor[rgb]{0.9,0.9,0.9}{
            \makecell[{{p{\linewidth}}}]{
                \texttt{\tiny{[GM$\langle$P2]}}
                \texttt{I need to analyze the referring expression and match it to one of the three grids.} \\
\\ 
\texttt{Let me break down the expression:} \\
\texttt{{-} "X's in the top{-}right corner and bottom{-}right area"} \\
\texttt{{-} "diagonal pattern of X's from top{-}left to bottom{-}right"} \\
\texttt{{-} "broken by filled cells in positions (1,2), (2,3), (3,4), and (4,5)" {-} using 1{-}indexed notation} \\
\texttt{{-} "empty cells at the four corners of the bottom{-}right 2x2 subgrid"} \\
\\ 
\texttt{Let me check each grid:} \\
\\ 
\texttt{**First Grid:**} \\
\texttt{```} \\
\texttt{□ □ X X □} \\
\texttt{X □ □ X X} \\
\texttt{□ X X □ X} \\
\texttt{X X □ □ X} \\
\texttt{□ X X □ □} \\
\texttt{```} \\
\texttt{{-} Top{-}right corner has X's at (1,3), (1,4)} \\
\texttt{{-} Bottom{-}right area: (3,5), (4,5) have X's} \\
\texttt{{-} Check diagonal from top{-}left: (1,1) empty, (2,2) empty, (3,3) X, (4,4) empty, (5,5) empty {-} not a clear diagonal} \\
\\ 
\texttt{**Second Grid:**} \\
\texttt{```} \\
\texttt{□ X □ X □} \\
\texttt{X □ X □ X} \\
\texttt{□ X □ X □} \\
\texttt{X □ X □ X} \\
\texttt{□ X □ X □} \\
\texttt{```} \\
\texttt{{-} This has a checkerboard pattern} \\
\texttt{{-} No X's in top{-}right corner specifically} \\
\texttt{{-} Diagonal from top{-}left: (1,1) empty, (2,2) empty, (3,3) empty, (4,4) empty, (5,5) empty {-} no diagonal pattern} \\
\\ 
\texttt{**Third Grid:**} \\
\texttt{```} \\
\texttt{□ X X □ X} \\
\texttt{X □ X X □} \\
\texttt{□ X □ X X} \\
\texttt{X □ X □ X} \\
\texttt{□ X X □ □} \\
\texttt{```} \\
\texttt{{-} Top{-}right corner: (1,3), (1,4) are X's, (1,5) is X ✓} \\
\texttt{{-} Bottom{-}right area: (3,4), (3,5), (4,5) are X's ✓} \\
\texttt{{-} Check diagonal pattern: (1,1) empty, (2,2) empty, (3,3) empty, (4,4) empty, (5,5) empty {-} main diagonal is empty} \\
\texttt{{-} But positions (1,2), (2,3), (3,4), (4,5) are filled: (1,2)=X ✓, (2,3)=X ✓, (3,4)=X ✓, (4,5)=X ✓} \\
\texttt{{-} Bottom{-}right 2x2 subgrid (rows 4{-}5, columns 4{-}5): (4,4)=empty, (4,5)=X, (5,4)=empty, (5,5)=empty {-} corners are (4,4), (4,5), (5,4), (5,5) {-} three empty, one filled. Wait, let me re{-}read: "empty cells at the four corners" {-} but (4,5) is X.} \\
\\ 
\texttt{Let me re{-}check: The bottom{-}right 2x2 subgrid corners would be positions (4,4), (4,5), (5,4), (5,5). In third grid: empty, X, empty, empty. That's not all empty.} \\
\\ 
\texttt{Wait {-} re{-}reading: "empty cells at the four corners of the bottom{-}right 2x2 subgrid" {-} this seems to mean the four corner positions of that subgrid should be empty. But (4,5) is X in third grid.} \\
\\ 
\texttt{Let me re{-}check first grid's bottom{-}right 2x2 (rows 4{-}5, cols 4{-}5): (4,4)=empty, (4,5)=X, (5,4)=empty, (5,5)=empty. Same issue.} \\
\\ 
\texttt{Actually re{-}reading more carefully {-} the expression says the diagonal is "broken by filled cells in positions (1,2), (2,3), (3,4), and (4,5)". This suggests these positions break the diagonal, meaning they should be filled. } \\
\\ 
\texttt{Checking third grid: (1,2)=X, (2,3)=X, (3,4)=X, (4,5)=X {-} all filled! ✓} \\
\\ 
\texttt{And "empty cells at the four corners of the bottom{-}right 2x2 subgrid" {-} looking at first grid: (4,4)=empty, (4,5)=X, (5,4)=empty, (5,5)=empty {-} not all empty.} \\
\\ 
\texttt{Hmm, let me re{-}interpret. Maybe "bottom{-}right 2x2 subgrid" refers to something else, or I need to check if the description matches third grid best overall.} \\
\\ 
\texttt{The third grid clearly has the filled cells at (1,2), (2,3), (3,4), (4,5) and has X's in top{-}right and bottom{-}right areas.} \\
\\ 
\texttt{Answer: third} \\
            }
        }
    }
     \\ \\

    \theutterance \stepcounter{utterance}  
    & & & \multicolumn{2}{p{0.3\linewidth}}{
        \cellcolor[rgb]{0.9,0.9,0.9}{
            \makecell[{{p{\linewidth}}}]{
                \texttt{\tiny{[GM$|$GM]}}
                \texttt{Invalid generated expression} \\
            }
        }
    }
    & & \\ \\

\end{supertabular}
}

\end{document}
